\section{Multivariate Division}

In the previous section, we established that the solvability of a system of polynomial equations reduces to the ideal membership problem, which states, given a polynomial $f$ and an ideal $I = \langle f_1, \dots, f_s \rangle$, is $f \in I$?

In the univariate case, where $R = k[x]$, this problem is trivialized by the Euclidean division algorithm. Since $k[x]$ is a principal ideal domain, any ideal is generated by a single greatest common divisor, say $g$. To check if $f \in \langle g \rangle$, one simply computes the remainder of $f$ divided by $g$. If the remainder is $0$, then $f \in I$; otherwise, it is not. Let's illustrate this with a simple example.

\begin{example}
\label{ex:univar_membership}
Consider the ideal $I = \langle x^2 - x \rangle$ in $\mathbb{Q}[x]$, and lets check if $h_1(x) = x^3 - x^2$ and $h_2(x) = x^3 - 1$ are elements of $I$. Dividing $h_1$ by $x^2 - x$ using polynomial long division 
$$
x^3 - x^2 = (x)(x^2 - x) + 0
$$
yields a remainder of $0$, confirming that $h_1 \in I$. Conversely, for $h_2$, we have
$$
x^3 - 1 = (x + 1)(x^2 - x) + (x - 1)
$$
with a non-zero remainder of $x - 1$, indicating that $h_2 \notin I$. This is a direct and trivial application of the division algorithm in one variable.
\end{example}

However, generalizing this logic to the multivariate ring $k[x_1, \dots, x_n]$ introduces complications. Consider another example.

\begin{example} 
\label{ex:multivar_membership}
Let $I = \langle x^2 y - x, x^2 + y^3 \rangle$ in the ring $\mathbb{Q}[x, y]$, and let's check if $h_3(x, y) = x^2y + y^4$ is an element of $I$. Here, $I$ is generated by two polynomials, its elements can be expressed as combinations of the generators: $a_1(x^2 y - x) + a_2(x^2 + y^3)$ for some $a_1, a_2 \in \mathbb{Q}[x, y]$. This example follows the same principle as before, but considering that $a_1$ and $a_2$ can be any polynomials in two variables, the search space is vastly larger. Other than impractical trial and error, there is no straightforward method to determine if such $a_1$ and $a_2$ exist that satisfy the equation $h_3 = a_1(x^2 y - x) + a_2(x^2 + y^3)$.
\end{example}

Our goal is to apply our solution for Example \ref{ex:univar_membership} to Example \ref{ex:multivar_membership} by developing a multivariate division algorithm. However, two main issues arise in this generalization.

Firstly, in Example \ref{ex:univar_membership}, there only exists one polynomial divisor to be used in the division algorithm, but in Example \ref{ex:multivar_membership}, there are two generators in the ideal, leading to multiple possible divisors to choose from at each step. We can overcome this problem by simply choosing one of the divisors in each step and keeping track of a quotient for each divisor. The final result in the univariate case is expressed as $f = aq + r$, where $a$ is the quotient, $q$ is the divisor, and $r$ is the remainder. In the multivariate case with multiple divisors, we can express the result as $f = a_1 q_1 + \dots + a_s q_s + r$, where each $a_i$ is the quotient corresponding to divisor $q_i$. This expression still captures the essence of the division algorithm, but now accommodates multiple divisors.

Secondly, notice that in the division algorithm we divide the leading term of the dividend by the leading term of the divisor, giving importance to what we consider the leading term, or the largest term in a polynomial. In the univariate case, this is straightforward since there is a natural ordering of terms by degree ($x^2 > x$). In the multivariate case, however, there is no inherent way to order terms like $x^2 y$ and $x y^2$. To address this, we need to follow a systematic rule for ordering monomials in multiple variables to ensure consistency in determining leading terms. We denote a monomial $x_1^{\alpha_1} \dots x_n^{\alpha_n}$ as $x^\alpha$, where $\alpha \in \mathbb{Z}_{\ge 0}^n$.

\begin{definition}[Monomial Ordering]
\label{def:monomial_ordering}
A monomial ordering on $k[x_1, \dots, x_n]$ is a relation $>$ on the set of monomials such that $>$ is a total ordering (any two monomials are comparable), $>$ is compatible with multiplication, as in, if $x^\alpha > x^\beta$, then $x^\alpha x^\gamma > x^\beta x^\gamma$ for all $\gamma$, and $>$ is a well-ordering (every non-empty set of monomials has a least element).
\end{definition}

The conditions of monomial ordering are set to guarantee consistency and termination of the division algorithm. The total ordering condition ensures that we can always identify a leading term in any polynomial. The compatibility condition guarantees that the ordering behaves predictably under multiplication, and the well-ordering condition is specially critical for computer science applications, as it guarantees that any algorithm reducing the size of a polynomial based on this order will eventually terminate. While there are infinitely many such orderings, two are very frequently used in computational algebra; lexicographic order (lex) and graded reverse lexicographic order (grevlex).

\begin{definition}[Lexicographic Order]
\label{def:lex_order}
Let $\alpha, \beta \in \mathbb{Z}_{\ge 0}^n$. The lexicographic order is defined by  $x^\alpha >_{lex} x^\beta$ if the leftmost non-zero entry of the vector difference $\alpha - \beta$ is positive.
\end{definition}
Intuitively, Lex behaves like a dictionary sort. For variables $x > y > z$, we have $x > y^5 z^3$ because the $x$-exponent dominates. On the other hand, $y^5 z^3 > y^4 z^{10}$ because the $y$-exponent is larger.

\begin{definition}[Graded Reverse Lexicographic Order]
\label{def:grevlex_order}
Graded reverse lexicographic order is defined by $x^\alpha >_{grevlex} x^\beta$ if $|\alpha| > |\beta|$, or if $|\alpha| = |\beta|$ and the rightmost non-zero entry of $\alpha - \beta$ is negative.
\end{definition}
Grevlex compares total degree first. For example, $x>y>z$, $y^2 > x z$ since $2 > 1$. If total degrees are equal, it favors monomials that have fewer variables from the end of the ordering. Thus, $x y > x z$ because the $y$-exponent is larger than the $z$-exponent. For computation, grevlex is often more efficient than lex as it tends to keep coefficients smaller.

It is worth noting that these definitions are specific instances of a more general framework. In fact, by Robbiano's Theorem \cite{Robbiano1986_TheoryOfGradedStructures}, any admissible monomial ordering on $k[x_1, \dots, x_n]$ can be represented by a non-singular matrix $M \in \text{GL}_n(\mathbb{Q})$. Given such a matrix $M$, we define the ordering $>_M$ by mapping the exponent vectors $\alpha$ and $\beta$ to the vector space $\mathbb{Q}^n$ via linear transformation and comparing the results lexicographically. Formally, we have
$$
x^\alpha >_M x^\beta \iff M\alpha >_{lex} M\beta
$$
where $>_{lex}$ on the right-hand side denotes the standard lexicographic comparison of vectors, essentially comparing the first non-zero entry of the difference.

To ensure that $>_M$ is a well-ordering compatible with multiplication, the matrix $M$ must satisfy specific positivity conditions, typically that the first non-zero entry of every column is positive. This matrix representation is particularly powerful because it allows us to construct custom orderings for specific applications, such as elimination theory, simply by designing the appropriate matrix.

\begin{example}
    For a Lex ordering with $x_1 > x_2 > \dots > x_n$, the corresponding matrix is the identity matrix $I_n$, as
    $$
    M_{lex} = \begin{pmatrix} 1 & 0 & \dots & 0 \\ 0 & 1 & \dots & 0 \\ \vdots & \vdots & \ddots & \vdots \\ 0 & 0 & \dots & 1 \end{pmatrix}
    $$
    comparing $M_{lex}\alpha$ with $M_{lex}\beta$ is identical to comparing $\alpha$ and $\beta$ directly.

    Grevlex can be represented by a matrix that first checks total degree (a row of 1s) and then applies negative weights to reverse the tie-breaking on variable indices. For $\mathbb{Q}[x, y, z]$, the matrix is 
    $$
    M_{grevlex} = \begin{pmatrix} 1 & 1 & 1 \\ 1 & 1 & 0 \\ 1 & 0 & 0 \end{pmatrix}
    $$
    where the first row compares the total degree $|\alpha|$ vs $|\beta|$. If they are equal, the subsequent rows enforce the reverse condition by prioritizing variables earlier in the sequence.
\end{example}

With a fixed monomial order, we can define the leading term of a polynomial $f$, denoted $\text{LT}(f)$, as well as its coefficient $\text{LC}(f)$ and monomial $\text{LM}(f)$.

Given a monomial order and a set of divisors $F = \{f_1, \dots, f_s\}$, we can formulate a multivariate polynomial long division algorithm with the goal to write $f$ as
$$ f = a_1 f_1 + \dots + a_s f_s + r $$
where no term in the remainder $r$ is divisible by any $\text{LT}(f_i)$. Algorithm \ref{alg:mult_polydivision} outlines this process.

\begin{definition}[Reduction]
\label{def:reduction}
When a polynomial $h$ is divided by a set of polynomials $F = \{f_1, \dots, f_s\}$ using the multivariate division algorithm, the resulting polynomial $r$ is written as $\text{reduce}(h, F)$ and is said that $h$ reduces to $r$.
\end{definition}

\begin{algorithm}[H]
\caption{Multivariate Division Algorithm} 
\label{alg:mult_polydivision}
\begin{algorithmic}
\Require $h, F=\{f_1, \dots, f_s\}, >$
\Ensure $r$
\State $p \gets h, r \gets 0$
\While{$p \neq 0$}
    \If{$\exists i$ s.t. $LT(f_i) \mid LT(p)$}
        \State $p \gets p - \frac{LT(p)}{LT(f_i)} f_i$
    \Else
        \State $r \gets r + LT(p), \enspace p \gets p - LT(p)$
    \EndIf
\EndWhile \ and \Return $r$
\end{algorithmic}
\end{algorithm}

We denote the result of this process as $r = \text{reduce}(h, F)$. If $\text{reduce}(h, F) = 0$, we have explicitly constructed coefficients $a_i$ such that $h = \sum a_i f_i$, proving that $h \in I$. However, unlike the univariate case, a non-zero remainder does \textit{not} prove that $h \notin I$. The output of the division algorithm depends heavily on the order in which the divisors $f_i$ are listed.

To demonstrate this, let us apply Algorithm \ref{alg:mult_polydivision} to Example \ref{ex:multivar_membership}, using the Lexicographic order with $x > y$. In this example, $f_1 = x^2 y - x$ and $f_2 = x^2 + y^3$, and we want to check if $h = x^2 y + y^4$ is in the ideal $I$. First, we list the divisors as $I = \langle f_1, f_2 \rangle$, so the algorithm checks $f_1$ first. In this case, the final result is $r = x + y^4 \neq 0$, suggesting that $h \notin I$. If we then list the divisors as $I = \langle f_2, f_1 \rangle$. The algorithm checks $f_2$ first, in which case the final result is $r = 0$, indicating that $h \in I$. Since we found a remainder of 0 in the second case, we have proven that $h \in I$. This example highlights that the standard multivariate division algorithm is not sufficient for the ideal membership problem because the remainder is not unique. Have in mind, that this is a simple example with only two generators, and in practice, ideals can have many generators and the orderings can be much more complex.

The failure of the division algorithm arises because the leading terms of the generators $\{f_1, \dots, f_s\}$ may not cover all the leading terms of the ideal $I$. Cancellation can occur between generators during linear combination, producing a new polynomial in $I$ with a leading term not divisible by any $\text{LT}(f_i)$. To resolve this, we could try to start the division process with a divisor that is an element of $I$ itself. However if the leading term of this element is not divisible by any $\text{LT}(f_i)$, we cannot proceed with the division. Thus, we need a generating set of $I$ such that every element of $I$ has its leading term divisible by some leading term of the generating set. A Gröbner basis is a generating set with this property.

\begin{definition}[Gröbner Basis]
\label{def:grobner_basis}
Let there be a fixed monomial order. A finite subset of generators $G = \{g_1, \dots, g_t\} \subset I$ is a Gröbner basis for a non-zero ideal $I$ if 
$$
\langle \text{LT}(g_1), \dots, \text{LT}(g_t) \rangle = \langle \{ \text{LT}(f) \mid f \in I \} \rangle
$$
\end{definition}

This definition implies two equivalent properties that resolve the issues discussed above:
\begin{center}
    \begin{varwidth}{\textwidth}
        \begin{enumerate}[label=(\roman*), font=\itshape]
            \item $f \in I \iff \text{reduce}(f, G) = 0$
            \item $\text{reduce}(f, G)$ is unique $\forall f \in R$ regardless of the order of $g_i$.
        \end{enumerate}
    \end{varwidth}
\end{center}

Property (i) directly answers Example \ref{ex:multivar_membership}, as if we could find a Gröbner basis $G$ for $I$, we could simply compute $\text{reduce}(h_3, G)$ to determine if $h_3 \in I$ without taking account of the order of generators. Property (ii) ensures that the multivariate division algorithm behaves consistently when using a Gröbner basis, making it a reliable tool for computations in polynomial ideals. This consistency and reliability makes Gröbner bases a widely used tool in computational algebraic geometry and computer algebra systems.